\documentclass[10pt,a4paper]{amsart}
\usepackage[margin=19mm]{geometry}
\usepackage{amsmath,amssymb,mathtools}
\usepackage[scr=rsfs,cal=euler]{mathalfa}
\usepackage{xfrac}
\usepackage[unicode,hidelinks,bookmarks=false]{hyperref}
\newcommand{\ie}{i.e.\ }
\newcommand{\cf}{cf.\ }
\newcommand{\resp}{respectively\ }
\newcommand{\Subsection}{\subsection}
\providecommand{\nobreakdash}{\nobreak\hbox{-}}
\setlength{\emergencystretch}{2em}
\multlinegap0pt
\usepackage{bm}
\usepackage{bbm}
\usepackage{dsfont}
\usepackage{xspace}
\usepackage{cancel}
\usepackage{mathtools}
\usepackage{amsthm}
\makeatletter
\@ifundefined{lem}{\newtheorem{lem}{Lemma}}{}
\@ifundefined{prop}{\newtheorem{prop}{Proposition}}{}
\makeatother
\newcommand{\R}{\mathbb{R}}
\title[Construction of 2-bubbles for the energy critical bi-harmoni]{Construction of 2-bubbles for the energy critical bi-harmonic Schr\"odinger equation}
\author{Jean-Baptiste Casteras, Ilkka Holopainen, L\'eonard Monsaingeon}
\date{Source: arXiv:2304.14744v3 (CC BY 4.0)}
\begin{document}
\maketitle

\noindent\emph{Source and licence.} Construction of 2-bubbles for the energy critical bi-harmonic Schr\"odinger equation. Version arXiv:2304.14744v3 (CC BY 4.0), distributed under the Creative Commons Attribution 4.0 International licence, as recorded in the arXiv licence metadata for this exact version and shown on the abstract page https://arxiv.org/abs/2304.14744v3. Full rights notice: licenses/arxiv-2304.14744-CC-BY-4.0.txt.\par\smallskip

\noindent\textit{Editorial scope.} Counted unit: the Prop whose source label is \texttt{4-4} in the article (source0.tex lines 1426--1451, source label \texttt{4-4}), delivered together with the article's own complete original proof of it (source0.tex lines 1453--1636, one \texttt{proof} environment of 184 lines). No mathematics is written, repaired or re-derived: statement and proof are the article's own text line for line. Required context retained uncounted: eq (lines 124--127); fe1 (lines 289--315); fe2 (lines 295--298); fe3 (lines 299--302); lem3.1 (lines 360--405); mode1 (lines 362--365); mode5 (lines 379--382); mode6 (lines 384--388); mode9 (lines 399--402); stcoe2 (lines 773--776); eqdeg (lines 808--812); propenere1 (lines 902--905); propenere2 (lines 907--911); propcoer111 (lines 1068--1086); lemcondinit (lines 1153--1172); condinit (lines 1156--1159); 4-3 (lines 1164--1166); eqg0 (lines 1167--1170); lemconsq (lines 1255--1266); lemproperA (lines 1343--1362); properA1 (lines 1348--1351); properA3 (lines 1357--1360). Every definition, display and label of those lines is retained unchanged. Omitted from this package and not redistributed: 1--123, 128--288, 303--359, 366--378, 383--383, 389--398, 403--772, 777--807, 813--901, 906--906, 912--1067, 1087--1152, 1160--1163, 1171--1254, 1267--1342, 1352--1356, 1361--1425, 1452--1452, 1637--2027. Nothing inside a retained excerpt is omitted or altered: every retained body is delivered exactly as the article's lines are written, so no omission receipt of any kind accompanies this package. Citations: the retained excerpts carry no citation command at all, so no bibliography entry of the article is transcribed and no reference is redistributed. Reference adaptation: none was needed and none was made; every reference inside the delivered unit resolves inside it, so no unresolved reference or citation is produced. Printed slips kept literal: no printed word, symbol, hypothesis or number of the counted statement or of its proof was altered. Layout and class: the article's own class is \texttt{amsart} with packages including amsmath, amsfonts, amssymb, amsthm, inputenc, fontenc, bbm, latexsym, xy, calc, multicol, hyperref, srcltx, xargs; the wrapper is the lane's pinned \texttt{amsart} editorial wrapper at 10pt on A4, which restates the theorem-like environments the retained text needs and adds the article's own macro definitions that the retained text needs (\texttt{\textbackslash R}), with extra wrapper packages (bm, bbm, dsfont, xspace, cancel, mathtools). The delivered numbering is the unit's own. Assets: no retained span contains a figure environment, a graphics-inclusion command, a PDF-page inclusion, a table or a TikZ picture; no image file is copied into this package, the package declares no asset dependency and no image file is redistributed with it. Licence: version arXiv:2304.14744v3 of this article is distributed under the Creative Commons Attribution 4.0 International licence, as recorded in the arXiv licence metadata for this exact version and shown on the abstract page https://arxiv.org/abs/2304.14744v3. A full rights notice is delivered at licenses/arxiv-2304.14744-CC-BY-4.0.txt. Characters: every retained excerpt is pure ASCII, so the delivered engine, XeLaTeX with its default Latin Modern text font, reports no missing character.
\begin{equation}
\label{eq}
i\partial_t u -\Delta^2 u + |u|^{2^\ast-2} u=0,\text{ in } \R^N.
\end{equation}

\begin{lem}
Let $N\geq 13$. For any $z_i \in \mathbb{C}$, $i=1,2,3$, we have
\begin{equation}
\label{fe1}
|f^\prime (z_1 +z_2)-f^\prime (z_1)|\lesssim |f^\prime (z_2)|,\ |f^\prime (z_1 +z_2)-f^\prime (z_1)|\lesssim |z_1|^{-\frac{N-12}{N-4}} |z_2|,\ if\ z_1 \neq 0,
\end{equation}
\begin{equation}
\label{fe2}
|f(z_1 +z_2) - f(z_1)|\lesssim |f^\prime (z_1)||z_2|+|f(z_2)|,
\end{equation}
\begin{equation}
\label{fe3}
|f(z_1 +z_2) - f(z_1)-f^\prime (z_1) z_2|\lesssim f(|z_2|),
\end{equation}
\begin{equation}
\label{fe4}
|f(z_1 +z_2) - f(z_1) -f^\prime (z_1) z_2|\lesssim |z_1|^{-\frac{N-12}{N-4}}|z_2|^2,
\end{equation}
\begin{equation}
\label{fe5}
|F(z_1 +z_2) - F(z_1)- \mathcal{R} (\overline{f(z_1)} z_2)|\lesssim |f^\prime (z_1)||z_2|^2+|F(z_2)|,
\end{equation}
\begin{equation}
\label{fe6}
|F(z_1 +z_2) - F(z_1) - \mathcal{R} (\overline{f(z_1)} z_2) - \mathcal{R} (\overline{f^\prime (z_1) z_2}z_2 ) |\lesssim |F(z_2)|.
\end{equation}
\end{lem}

\begin{prop}\label{lem3.1}
Let $c>0$ be an arbitrarily small constant. Let $T_0=T_0(c)<0$ with $|T_0|$ large enough and $T<T_1 \leq T_0$. Suppose that for $T\leq t \leq T_1$, we have
\begin{equation}
\label{mode1}
|\zeta (t) +\frac{\pi}{2}| \leq |t|^{-\frac{3}{N-12}},
\end{equation}
\begin{equation}
\label{mode2}
|\mu (t) -1| \leq |t|^{-\frac{3}{N-12}},
\end{equation}
\begin{equation}
\label{mode3}
|\theta (t)| \leq |t|^{-\frac{1}{N-12}},
\end{equation}
\begin{equation}
\label{mode4}
|\lambda (t) - \tilde{C} |t|^{-\frac{2}{N-12}} |\leq 
|t|^{-\frac{5}{2(N-12)}},
\end{equation}
\begin{equation}
\label{mode5}
\|g\|_{\mathcal{E}} \leq |t|^{-\frac{N-3}{2(N-12)}}.
\end{equation}
Then, for $T\leq t\leq T_1$, there hold
\begin{equation}
\label{mode6}
|\zeta^\prime (t)|\leq c |t|^{-\frac{N-5}{N-12}},
%|t|^{-\frac{N-9}{N-12}},
\end{equation}
\begin{equation}
\label{mode7}
|\mu^\prime (t)| \leq c |t|^{-\frac{N-5}{N-12}},
%|t|^{-\frac{N-9}{N-12}},
\end{equation}
\begin{equation}
\label{mode8}
\left|\lambda^\prime (t) - \frac{C_1}{2\|W\|_{L^2}^2} \lambda(t)^{\frac{N-10}{2}}\right|\leq c |t|^{-\frac{2N-19}{2(N-12)}} ,
\end{equation}
and
\begin{equation}
\label{mode9}
\left|\theta^\prime (t) + \frac{C_2}{2\|W\|_{L^2}^2}\theta(t)\lambda(t)^{\frac{N-12}{2}} -\frac{K(t)}{2\lambda(t)^4 \|W\|_{L^2}^2} \right|\leq c |t|^{-\frac{N-11}{N-12}} ,
\end{equation}
where
$$K= - \langle  e^{i\theta}  \Lambda W_\lambda ,  f(e^{i\zeta} W_\mu +e^{i\theta} W_\lambda +g) - f(e^{i\zeta} W_\mu +e^{i\theta} W_\lambda) -f^\prime (e^{i\theta} W_\lambda +e^{i\zeta} W_\mu) g \rangle .$$
\end{prop}

\begin{equation}
\label{stcoe2}
\left|\frac{d}{dt}a_1^- (t) + \frac{\nu}{\mu (t)^4} a_1^- (t)\right| \leq \frac{c}{\mu(t)^4} |t|^{-\frac{N}{2(N-12)}},
\end{equation}

\begin{align}
\label{eqdeg}
\partial_t g &=-i \Delta^2 g +i\big(f(e^{i\zeta} W_\mu +e^{i\theta} W_\lambda +g) - f(e^{i\zeta} W_\mu ) - f(e^{i\theta }W_\lambda)\big)\\
&-  \zeta^\prime i e^{i\zeta} W_\mu + \frac{\mu^\prime}{\mu} e^{i\zeta}\Lambda W_\mu -\theta^\prime i e^{i\theta} W_\lambda + \frac{\lambda^\prime}{\lambda}e^{i\theta} \Lambda W_\lambda  .\nonumber
\end{align}

\begin{equation}
\label{propenere1}
|E(u)-2E(W)|\leq C \left((|\zeta +\pi /2| +|\mu -1| +|\theta| +\lambda )\lambda^{\frac{N-4}{2}} +\|g\|_{\mathcal{E}}^2 \right), 
\end{equation}

\begin{align}
\label{propenere2}
\|g\|_{\mathcal{E}}^2 &+C_0 \theta \lambda^{\frac{N-4}{2}} \leq C \Big(\lambda^{(N-4)/2} \big(|\zeta +\pi /2| +|\mu -1|+|\theta|^3 +\lambda\big) \nonumber\\
&+E(u) - 2E(W) +\sum_{j=1}^2 \big((a_j^+)^2 + (a_j^- )^2 \big)\Big).
\end{align}

\begin{prop}
\label{propcoer111}
There exist constants $c,C>0$ such that for any $\theta \in \R$, $\lambda >0$ and for any complex valued radial $g\in \mathcal{E}$, there holds
\begin{align*}
&\int_{\R^N} |\Delta g|^2 dx - \mathcal{R} \int_{\R^N} \bar{g} f^\prime (e^{i\theta} W_\lambda) g dx\\
& \geq c \int_{\R^N} |\Delta g|^2 dx - C (\langle \lambda^{-2} e^{i\theta} W_\lambda , g\rangle^2 + \langle \lambda^{-2} i e^{i\theta} \Lambda W_\lambda ,g \rangle^2 +\langle \alpha_{\theta ,\lambda}^+ ,g \rangle^2 +\langle \alpha_{\theta ,\lambda}^- ,g\rangle^2)
\end{align*}
If $r_1 >0$ is large enough, then we have
\begin{align*}
&(1-2c)\int_{|x|\leq r_1} |\Delta g|^2 dx+c \int_{|x_1| \geq r_1} |\Delta g|^2 dx - \mathcal{R} \int_{\R^N} \bar{g} f^\prime (e^{i\theta} W_\lambda) g dx\\
& \geq c \int_{\R^N} |\Delta g|^2 dx - C (\langle \lambda^{-2} e^{i\theta} W_\lambda , g\rangle^2 + \langle \lambda^{-2} i e^{i\theta} \Lambda W_\lambda ,g \rangle^2 +\langle \alpha_{\theta ,\lambda}^+ ,g \rangle^2 +\langle \alpha_{\theta ,\lambda}^- ,g\rangle^2)
\end{align*}
If $r_2>0$ is small enough, we have
\begin{align*}
&(1-2c)\int_{|x|\geq r_2} |\Delta g|^2 dx+c\int_{|x|\leq r_2} |\Delta g|^2 dx - \mathcal{R} \int_{\R^N} \bar{g} f^\prime (e^{i\theta} W_\lambda) g dx\\
& \geq c \int_{\R^N} |\Delta g|^2 dx - C (\langle \lambda^{-2} e^{i\theta} W_\lambda , g\rangle^2 + \langle \lambda^{-2} i e^{i\theta} \Lambda W_\lambda ,g \rangle^2 +\langle \alpha_{\theta ,\lambda}^+ ,g \rangle^2 +\langle \alpha_{\theta ,\lambda}^- ,g\rangle^2)
\end{align*}

\end{prop}

\begin{lem}
\label{lemcondinit}
There exists $T_0<0$ such that for all $T\leq T_0$ and for all $\lambda^0$, $a_1^0$, $a_2^0$ satisfying
\begin{equation}
\label{condinit}
|\lambda^0 - \tilde{C} |T|^{-\frac{2}{N-12}}|\leq \frac{1}{2} |T|^{-\frac{5}{2(N-12)}},\ |a_1^0|\leq \frac{1}{2} |T|^{-\frac{N}{2(N-12)}},\ |a_2^0|\leq \frac{1}{2} |T|^{-\frac{N}{2(N-12)}},
\end{equation}
there exists $g^0 \in X^2 = \dot{H}^3 (\R)^N \cap \dot{H}^2 (\R^N)$ satisfying
\begin{equation}\label{4-2}
\langle \Lambda W , g^0\rangle = \langle iW , g^0\rangle = \langle i \Lambda W_{\lambda^0},g^0 \rangle = \langle - W_{\lambda^0}, g^0\rangle =0,
\end{equation}
\begin{equation}\label{4-3}
\langle \alpha^-_{-\frac{\pi}{2},1} , g^0 \rangle =0,\ \langle \alpha^+_{-\frac{\pi}{2},1}, g^0 \rangle = a_1^0 ,\ \langle \alpha_{0,\lambda^0}^- , g^0 \rangle =0 ,\ \langle \alpha_{0,\lambda^0}^+,g^0 \rangle = a_2^0,
\end{equation}
\begin{equation}
\label{eqg0}
\|g^0\|_{\mathcal{E}} \lesssim |T|^{-\frac{N}{2(N-12)}} .
\end{equation}
Moreover, $g^0$ is continuous in the $X^2$ topology with respect to the parameters $\lambda^0$, $a_1^0$ and $a_2^0$.
\end{lem}

\begin{lem}
\label{lemconsq}
For any $c>0$ and $R>0$, there exists a radial function $q=q(c,R) \in C^{5,1} (\R^N)$ satisfying the following properties:
\begin{itemize}
\item[(1)] $q(x)= \frac{1}{2} |x|^2$ for $|x|\leq R$.
\item[(2)] There exists $\tilde{R}>0$ (depending on $c$ and $R$) such that $q(x) \equiv \text{Const}$ for $|x|\geq \tilde{R}$.
\item[(3)] $|\nabla q (x)|\lesssim |x|$ and $|\Delta q(x) | \lesssim 1 $ for all $x\in \R^N$, with constants independent of $c$ and $R$.
\item[(4)] $\partial_{rr}^2 q(x) \geq - c$ for all $x\in \R^N$. %$\sum_{1\leq j,k \leq N} (\partial_{x_j x_k}^2 q(x)) \overline{v_j} v_k \geq - c \sum_{j=1}^N |v_j|^2$ for all $x\in \R^N$, $v_j \in \mathbb{C}$.
\item[(5)] $\big(2\partial_{rr}^2(\Delta q)+\Delta^2 q\big)(x) \leq c |x|^{-2}$ and $-\Delta^3 q (x) \leq c |x|^{-4} $  for all $x\in \R^N$.
% and $\lim_{|x|\rightarrow R^+} \partial_r \Delta^2 q^2 (x)>0.$
\end{itemize}
\end{lem}

\begin{lem}
\label{lemproperA}
\begin{itemize}
\item[(i)] For $\lambda>0$, the families $\{A(\lambda )\}, \{A_0 (\lambda) \}, \{\lambda \partial_\lambda A(\lambda)\}$, and $\{\lambda \partial_\lambda A_0 (\lambda)\}$ are bounded in $\mathcal{L} (\mathcal{E} ; \dot{H}^{-2})$ and the families $\{\lambda A(\lambda)\}$ and  $\{\lambda A_0 (\lambda)\}$ are bounded in $\mathcal{L} (\mathcal{E} ; L^2)$.
\item[(ii)] For any complex-valued function $h_1 ,h_2 \in X^2 (\R^N)$ and $\lambda>0$, we have
\begin{equation}
\label{properA1}
\langle A(\lambda)h_1 , f(h_1 +h_2) - f(h_1) - f^\prime (h_1) h_2\rangle = - \langle A(\lambda)h_2 , f(h_1 +h_2) - f(h_1) \rangle ,
\end{equation}
\begin{equation}
\label{properA2}
\langle h_1 ,A_0 (\lambda) h_2 \rangle =- \langle A_0 (\lambda ) h_1 ,h_2 \rangle .
\end{equation}
\item[(iii)] For any $c_0>0$, if we choose $c$ in Lemma \ref{lemconsq} small enough, then for all $h\in X^2$,
\begin{equation}
\label{properA3}
\langle A_0 (\lambda )h, -\Delta^2 h \rangle \leq \dfrac{c_0}{\lambda^4 } \|h\|^2_{\mathcal{E}} - \dfrac{1}{\lambda^4} \int_{|x|\leq R \lambda} |\Delta h(x)|^2 dx .
\end{equation}
\end{itemize}
\end{lem}

\begin{prop}\label{4-4}
There exists $T_0<0$ with the following property. Let $T<T_1<T_0$ and let $\lambda^0$, $a_1^0$ and $a_2^0$ satisfy \eqref{condinit}. Let $g^0 \in X^2$ be given by Lemma \ref{lemcondinit} and consider the solution $u(t)$ to \eqref{eq} with initial data $u(T)=-iW+W_{\lambda_0}+g^0$. Suppose that $u(t)$ exists on $[T,T_1]$ and that \eqref{mode1}- \eqref{mode5} hold for any $t\in [T,T_1]$ as well as
\begin{equation}
\label{lastcondai}
|a_1^+ (t)| \leq |t|^{-\frac{N}{2(N-12)}},\ |a_2^+ (t)|\leq |t|^{-\frac{N}{2(N-12)}} .
\end{equation}
Then, for any $t\in [T,T_1]$, we have
\begin{equation}
\label{lastconde1}
\big|\zeta (t) +\frac{\pi}{2}\big| \leq \frac{1}{2} |t|^{-\frac{3}{N-12}},
\end{equation}
\begin{equation}
\label{lastconde2}
|\mu (t) -1|\leq \frac{1}{2} |t|^{-\frac{3}{N-12}},
\end{equation}
\begin{equation}
\label{lastconde3}
|\theta (t)|\leq \frac{1}{2} |t|^{-\frac{1}{N-12}},
\end{equation}
and
\begin{equation}
\label{lastconde4}
\|g(t)\|_{\mathcal{E}} \leq \frac{1}{2}|t|^{- \frac{N-3}{2(N-12)}}.
\end{equation}

\end{prop} 

\begin{proof}
Integrating \eqref{mode6} over $[T,t]$ with the initial value $\zeta (T)= -\frac{\pi}{2}$ we obtain
\begin{align*}
\left|\zeta (t)+\dfrac{\pi}{2}\right|&= |\zeta (t) - \zeta (T)| =
\left|\int_T^t \zeta^\prime (s) ds\right|\leq c\int_T^t  |s|^{-\frac{N-5}{N-12}} ds \le
 \frac{c(N-12)}{7}|t|^{-\frac{7}{N-12}}\\
& \le\frac{1}{2}|t|^{-\frac{3}{N-12}}. 
\end{align*}
% if $c\le\tfrac{7}{2(N-12)}$. 
The estimate \eqref{lastconde2} follows similarly.

Next we show that
\begin{equation}\label{a1_a2_}
|a_1^-(t)|<|t|^{-\frac{N}{2(N-12)}}\quad\text{and}\quad
|a_1^-(t)|<|t|^{-\frac{N}{2(N-12)}}
\end{equation}
for $t\in [T,T_1]$. The claims hold for $t=T$ by \eqref{4-3}.
Let $T_2\in (T,T_1)$ be the last time for which \eqref{a1_a2_} holds for $t\in [T,T_2)$. Suppose, for example, that $a_1^-=|T_2|^{-\frac{N}{2(N-12)}}$. 
Then \eqref{stcoe2} implies that $\frac{d}{dt}a_1^-(T_2)<0$ provided $c<\nu$. This contradicts the assumption that $a_1^-<|T_2|^{-\frac{N}{2(N-12)}}$
for $t<T_2$ (note that $T_2<T_0<0$). The other inequality in \eqref{a1_a2_} follows similarly.

Next we will prove that, for all $c_0>0$,
\begin{equation}\label{(4-18)}
|\theta(t)|\le c_0|t|^{-\frac{1}{N-12}}
\end{equation}
for $t\in [T,T_1]$ provided $|T_0|$ is chosen large enough depending on $c_0$.
By \eqref{propenere1}, \eqref{condinit}, \eqref{eqg0} and the conservation of the energy, we have
\begin{align*}
|E(u)- &2E(W)|=|E(u(T))-2E(W)|\\
&\lesssim \big(|\zeta (T) +\pi /2| +|\mu (T) -1| +|\theta (T)| +\lambda (T) \big)\lambda (T)^{\frac{N-4}{2}} +\|g(T)\|_{\mathcal{E}}^2\\
&=(\lambda^0)^{\frac{N-2}{2}}+\|g^0\|_{\mathcal{E}}^2\\
&\lesssim |T|^{-\frac{N-2}{N-12}} \le |t|^{-\frac{N-2}{N-12}}.
\end{align*}
So, using \eqref{propenere2}, we get
\begin{align*}
\theta \lambda^{\frac{N-4}{2}}& \lesssim \big(|\zeta +\pi/2|+|\mu -1| +|\theta|^3 +\lambda\big) \lambda^{\frac{N-4}{2}}
+|t|^{-\frac{N-2}{N-12}}+|t|^{-\frac{N}{N-12}}\\
&\lesssim |t|^{-\frac{N-2}{N-12}}.
\end{align*}
We deduce that
\begin{equation}\label{4-19}
\theta(t) \lesssim |t|^{-\frac{2}{N-12}}
\ll |t|^{-\frac{1}{N-12}}.
\end{equation}
In order to prove a converse estimate
\begin{equation}\label{4-20}
\theta(t)\ge -c_0|t|^{-\frac{1}{N-12}}
\end{equation}
we define
\[
\psi(t):=\theta (t) - \dfrac{1}{4\|W\|_{L^2}^2} \langle g(t), i A_0 (\lambda (t)) g(t)\rangle.
\]
We claim that 
\begin{equation}\label{4-21}
\psi^\prime (t)\ge -c_1|t|^{-\frac{N-11}{N-12}}
\end{equation}
for $t\in [T,T_1]$ with the constant $c_1>0$ as small as we wish by taking $|T_0|$ large enough.
Taking $|T_0|$ large enough and choosing $c=c_1/4$ in Proposition~\ref{lem3.1} we obtain, by using \eqref{mode9} and \eqref{4-19}, 
\begin{align}\label{4-22}
\psi^\prime &\geq 
-\frac{C_2}{2\|W\|_{L^2}^2}\theta\lambda^{\frac{N-12}{2}}
+\frac{K}{2\lambda^4 \|W\|_{L^2}^2} -\frac{c_1}{4}|t|^{-\frac{N-11}{N-12}} -\frac{1}{4\|W\|_{L^2}^2}\frac{d}{dt} \langle g, i A_0 (\lambda) g\rangle \nonumber\\
&\geq \frac{1}{2\|W\|_{L^2}^2} \left(\frac{K}{\lambda^4} - \frac{1}{2} \frac{d}{dt} \langle g, iA_0 (\lambda ) g \rangle \right) - \frac{c_1}{2}|t|^{-\frac{N-11}{N-12}}.
\end{align}
Hence we need to compute $ \frac{1}{2} \frac{d}{dt} \langle g, iA_0 (\lambda ) g \rangle$ up to terms of order $\ll |t|^{-\frac{N-11}{N-12}}$.
Using that $iA_0 (\lambda)$ is symmetric, we find
\begin{equation}\label{4-23}
\frac{1}{2} \frac{d}{dt} \langle g, iA_0 (\lambda ) g \rangle  = \frac{1}{2}\lambda^\prime \langle g, i\partial_\lambda A_0 (\lambda) g \rangle + \langle \partial_t g , iA_0 (\lambda) g\rangle.
\end{equation}
Since $\|i\lambda \partial_\lambda A_0 (\lambda) g \|_{\dot{H}^{-2}} \lesssim \|g\|_{\mathcal{E}}$, we get that
$$|\lambda^\prime \langle g, i\partial_\lambda A_0 (\lambda) g \rangle| \lesssim \big|\frac{\lambda^\prime}{\lambda}\big| \|g\|_{\mathcal{E}}^2
\ll |t|^{-\frac{N-11}{N-12}} .$$
To estimate the second term in \eqref{4-23} we write
 $\partial_t g$ as in \eqref{eqdeg}. First we estimate 
 $$C_3:=\langle -  \zeta^\prime i e^{i\zeta} W_\mu + \frac{\mu^\prime}{\mu} e^{i\zeta}\Lambda W_\mu -\theta^\prime i e^{i\theta} W_\lambda + \frac{\lambda^\prime}{\lambda} \Lambda W_\lambda , iA_0(\lambda) g \rangle . $$
Since $|\zeta^\prime (t)| +|\frac{\mu^\prime (t)}{\mu (t)}| + |\theta^\prime (t)|+ |\frac{\lambda^\prime (t)}{\lambda (t)}|\lesssim |t|^{-1} $, and $\|A_0 (\lambda) g\|_{\dot{H}^{-2}}\lesssim \|g\|_{\mathcal{E}}$, we get
$$|C_3|\lesssim |t|^{-1}\|g\|_{\mathcal{E}}
\le |t|^{-\frac{3N-27}{2(N-12)}}\ll |t|^{-\frac{N-11}{N-12}}. $$
\iffalse
So overall, we are left with estimating
$$\dfrac{1}{2} \dfrac{d}{dt} \langle g, iA_0 (\lambda ) g \rangle \approx \langle    -\Delta^2 g + f(e^{i\zeta} W_\mu +e^{i\theta}W_\lambda +g)- f(e^{i\zeta} W_\mu) - f(e^{i\theta} W_\lambda)  , A_0 (\lambda )g\rangle .$$
\fi
Hence we have
\begin{align*}
&
\frac{1}{2} \frac{d}{dt} \langle g, iA_0 (\lambda ) g \rangle \\
&\quad= \langle -\Delta^2 g + f(e^{i\zeta} W_\mu +e^{i\theta}W_\lambda +g)- f(e^{i\zeta} W_\mu) - f(e^{i\theta} W_\lambda)  , A_0 (\lambda )g\rangle\\
&\qquad
+o\big(|t|^{-\frac{N-11}{N-12}}\big) .
\end{align*}
Next we notice that the function $A_0 (\lambda ) g$ is supported in the ball of radius $\tilde{R}\lambda$ and in this ball, we have $W_\lambda \ll W_\mu$. Consequently,  
$|f(e^{i\zeta} W_\mu +e^{i\theta}W_\lambda) - f(e^{i\zeta}W_\mu)- f (e^{i\theta}W_\lambda)|\lesssim |W|_\lambda^{\frac{8}{N-4}}$ by Taylor's expansion \eqref{fe2}. On the other hand, we have
$$\|W_\lambda^{\frac{8}{N-4}}\|_{L^2 (|x|\leq \tilde{R} \lambda)}=\lambda^{\frac{N-4}{2}} \|W^{\frac{8}{N-4}}\|_{L^2 (|x|\leq \tilde{R})}\lesssim \lambda^{\frac{N-4}{2}}\lesssim |t|^{-\frac{N-4}{N-12}},$$
and, by Lemma \ref{lemproperA} (i),
 $\|A_0 (\lambda) g\|_{L^2} \lesssim \lambda^{-1}\|g\|_{\mathcal{E}}.$
So, we deduce from the three previous estimates and the Cauchy-Schwarz inequality that
\begin{align*}
\frac{1}{2} \frac{d}{dt} \langle g, i A_0 (\lambda) g\rangle & = \langle - \Delta^2 g + f(e^{i\zeta}W_\mu +e^{i\theta} W_\lambda +g   )-   f(e^{i\zeta}W_\mu +e^{i\theta} W_\lambda ) , A_0 (\lambda )g \rangle\\
&\qquad
+o\big(|t|^{-\frac{N-11}{N-12}}\big) .
\end{align*}

Using \eqref{properA1}, \eqref{properA3} and since $A_0 (\lambda) g = \frac{2}{N\lambda^4}\Delta q \big(\frac{\cdot}{\lambda}\big) g +A(\lambda ) g$,
\begin{align*}
&\frac{d}{dt} \big\langle g, i A_0 (\lambda) g\big\rangle 
\leq \frac{2c_0}{\lambda^4} \|g\|_{\mathcal{E}}^2\\
&- \frac{1}{\lambda^4} \Big(\int_{|x|\leq R\lambda} |\Delta g|^2 dx -\big\langle f(e^{i\zeta}W_\mu +e^{i\theta} W_\lambda +g) - f(e^{i\zeta} W_\mu +e^{i\theta} W_\lambda), \frac{2}{N} \Delta q \big(\frac{\cdot}{\lambda} \big) g\big\rangle  \Big) \\
&-\big\langle A(\lambda) (e^{i\zeta} W_\mu +e^{i\theta} W_\lambda), f (e^{i\zeta} W_\mu +e^{i\theta} W_\lambda+g) - f(e^{i\zeta} W_\mu +e^{i\theta} W_\lambda)\\
&\qquad - f^\prime (e^{i\zeta} W_\mu +e^{i\theta} W_\lambda)g\big\rangle .
\end{align*}
By applying Taylor's expansion \eqref{fe3}, we get 
\begin{align*}
& \| f (e^{i\zeta}W_\mu +e^{i\theta} W_\lambda +g) - f (e^{i\zeta} W_\mu +e^{i\theta} W_\lambda )  - f^\prime (e^{i\zeta} W_\mu +e^{i\theta} W_\lambda) g \|_{L^{\frac{2N}{N+4}}}\\
& \lesssim \|f(g) \|_{L^{\frac{2N}{N+4}}}  \lesssim  \|g\|_{\mathcal{E}}^{\frac{N+4}{N-4}} \ll \|g \|_{\mathcal{E}}.
\end{align*}
From \eqref{fe1}, we also have
\begin{align*}
\|f^\prime (e^{i\zeta} W_\mu +e^{i\theta} W_\lambda) - f^\prime (e^{i\theta} W_\lambda))g \|_{L^{\frac{2N}{N+4}}(|x| \leq \tilde{R} \lambda) } & \lesssim \|f^\prime (e^{i\zeta} W_\mu)\|_{L^{\frac{N}{4}} (|x|\leq \tilde{R} \lambda)} \|g\|_{\mathcal{E}}\\
&\ll \|g\|_{\mathcal{E}}.
\end{align*}
Combining the two previous estimate and since, by Sobolev's embedding, we have $\|\frac{1}{N} \Delta q (\frac{.}{\lambda}) g \|_{L^{\frac{2N}{N-4}}} \lesssim \|g\|_{\mathcal{E}}$, we obtain
\begin{align*}
&\big|\big\langle f(e^{i\zeta}W_\mu +e^{i\theta} W_\lambda +g) - f(e^{i\zeta}W_\mu +e^{i\theta} W_\lambda) , \frac{1}{N}\Delta q \big(\frac{\cdot}{\lambda}\big) g  \big\rangle - \big\langle f^\prime (e^{i\theta}W_\lambda)g , \frac{1}{N}\Delta q \big(\frac{\cdot}{\lambda}\big) g \big\rangle\big|\\
& \ll %\|g\|_{\mathcal{E}}^{1+\frac{N+4}{N-4}}+
\|g\|_{\mathcal{E}}^2 . 
\end{align*}

Since $\frac{1}{N}\Delta q \big(\frac{x}{\lambda}\big)=1$, for $|x|\leq R \lambda$, and $\|f^\prime (e^{i\theta} W_\lambda)\|_{L^{\frac{N}{4}} (|x|\geq R\lambda)} \ll 1$ for $R$ large enough, we have
$$\big|\big\langle f(e^{i\zeta}W_\mu +e^{i\theta} W_\lambda +g) - f(e^{i\zeta}W_\mu +e^{i\theta} W_\lambda) , \frac{1}{N}\Delta q \big(\frac{\cdot}{\lambda}\big) g  \big\rangle - \langle f^\prime (e^{i\theta}W_\lambda)g ,  g \rangle\big| \ll  |t|^{-\frac{N-3}{N-12}}.$$
%\|g\|_{\mathcal{E}}^2 .$$

By Proposition \ref{propcoer111}, \eqref{a1_a2_}, and the assumption \eqref{lastcondai}, we have, for some constant $c_3>0$ that can be chosen as small as we wish by enlarging $R$, that
$$\int_{|x|\leq R\lambda} |\Delta g|^2 dx - \langle f^\prime (e^{i\theta } W_\lambda) g ,g\rangle \geq - c_3 \|g\|_{\mathcal{E}}^2 $$
 up to some negligible terms.
Hence
\begin{align*}
\frac{2c_0}{\lambda^4} \|g\|_{\mathcal{E}}^2 & - \frac{1}{\lambda^4} \Big(\int_{|x|\leq R\lambda} |\Delta g|^2 dx \\
&\qquad -\big\langle f (e^{i\zeta}W_\mu +e^{i\theta} W_\lambda +g) - f (e^{i\zeta} W_\mu +e^{i\theta} W_\lambda ) , \frac{1}{N} \Delta q \big(\frac{\cdot}{\lambda} \big) g\big\rangle  \Big) \\
& \leq \frac{c_2}{\lambda^4} \|g\|_{\mathcal{E}}^2
\end{align*}
with a positive constant $c_2$ that can be taken as small as we wish by enlarging $R$.

Since the support of $A(\lambda)(e^{i\zeta} W_\mu)$ is contained in $|x|\leq \tilde{R} \lambda$ and 
\iffalse
$\|A(\lambda)(e^{i\zeta} W_\mu)\|_{L^\infty}\lesssim \lambda^{-2}$, we have
$$\|A(\lambda)(e^{i\zeta} W_\mu)\|_{L^{\frac{2N}{N-4}}} \lesssim (\lambda^N \lambda^{-\frac{4N}{N-4} } )^{\frac{N-4}{2N}}=\lambda^{\frac{N-4}{2}-2}.  $$
\fi
$\|A(\lambda)(e^{i\zeta} W_\mu)\|_{L^\infty}\lesssim \lambda^{-4}$, so we have
$$\|A(\lambda)(e^{i\zeta} W_\mu)\|_{L^{\frac{2N}{N-4}}} \lesssim (\lambda^N \lambda^{-\frac{4N}{N-4} } )^{\frac{N-4}{2N}}=\lambda^{\frac{N-12}{2}}\approx |t|^{-1}.  $$

On the other hand, by \eqref{fe3}, we get
$$\|f(e^{i\zeta}W_\mu +e^{i\theta} W_\lambda +g) - f (e^{i\zeta}W_\mu +e^{i\theta} W_\lambda) - f^\prime (e^{i\zeta} W_\mu +e^{i\theta} W_\lambda) g \|_{L^{\frac{2N}{N+4}}} \lesssim \|g\|_{\mathcal{E}}^{\frac{N+4}{N-4}}. $$

Observe that
$$|A(\lambda)(e^{i\theta} W_\lambda ) - \frac{1}{\lambda^4} e^{i\theta} \Lambda W_\lambda| \lesssim \frac{W_\lambda}{\lambda^4},$$
and $A(\lambda ) W = \frac{1}{\lambda^4} \Lambda W_\lambda$ for $|x|\leq R \lambda$, so we obtain
\begin{align*}
\big|\big\langle A(\lambda) (e^{i\theta }W_\lambda)- \frac{1}{\lambda^4} e^{i\theta}\Lambda W_\lambda,&\, f(e^{i\zeta}W_\mu +e^{i\theta} W_\lambda +g)- f (e^{i\zeta} W_\mu +e^{i\theta} W_\lambda)\\
&\, - f^\prime (e^{i\zeta} W_\mu +e^{i\theta} W_\lambda)g\big\rangle\big|\\
 \lesssim \frac{1}{\lambda^4} \int_{|x|\geq R\lambda} W_\lambda\big|f(e^{i\zeta}W_\mu &+e^{i\theta} W_\lambda +g)- f (e^{i\zeta} W_\mu +e^{i\theta} W_\lambda)\\
&\, - f^\prime (e^{i\zeta} W_\mu +e^{i\theta} W_\lambda)g \big| dx.
\end{align*}
Since $|\zeta -\theta|\approx \frac{\pi}{2} $, we have $|e^{i\zeta}W_\mu + e^{i\theta}W_\lambda | \gtrsim W_\lambda $, and therefore \eqref{fe3} implies that
$$W_\lambda \big|f(e^{i\zeta}W_\mu +e^{i\theta} W_\lambda +g)- f (e^{i\zeta} W_\mu +e^{i\theta} W_\lambda) - f^\prime (e^{i\zeta} W_\mu +e^{i\theta} W_\lambda)g\big| \lesssim W_\lambda^{\frac{8}{N-4}} |g|^2 .$$
%Computing the last term by integrated over $|x|\geq R\lambda$, we find overall that
Integrating the last term over the region $|x|\geq R\lambda$, we find overall that
$$\frac{1}{2} \dfrac{d}{dt} \langle g ,iA_0 (\lambda ) g \rangle \leq \frac{c_1}{2\lambda^4 } \|g\|_{\mathcal{E}}^2 +\dfrac{K}{\lambda^4}.$$
So we proved that 
\[
\psi^\prime (t) \geq - c_1 |t|^{-\frac{N-11}{N-12}}.
\]
Since $\theta (T)=0 $, we have $|\psi (T)|\lesssim \|g\|_{\mathcal{E}}^2\ll |t|^{-\frac{1}{N-12}} $. Integration of \eqref{4-21} over $[T,t]$ yields $\psi(t) \gtrsim -c_1|t|^{-\frac{1}{N-1}}$. But, on the other hand,
$$|\langle g(t) , A_0 (\lambda ) g(t) \rangle| \lesssim \|g(t)\|_{\mathcal{E}}^2\le |t|^{-\frac{N-3}{N-12}}
\ll |t|^{-\frac{1}{N-12}}.$$
Hence $\theta (t) \gtrsim -c_1|t|^{-\frac{1}{N-12}}$ that implies \eqref{4-20} provided $c_1>0$ is chosen small enough. We have proved \eqref{lastconde3}.

Finally, from \eqref{propenere2} we get
$\|g\|_{\mathcal{E}}^2 +C_0 \theta \lambda^{\frac{N-4}{2}}\le C|t|^{-\frac{N-2}{N-12}}$, and therefore
\[
\|g\|_{\mathcal{E}}^2 \le C_0 \theta \lambda^{\frac{N-4}{2}} +C|t|^{-\frac{N-2}{N-12}}\le \frac{1}{8}|t|^{-\frac{N-3}{N-12}} +C|t|^{-\frac{N-2}{N-12}}
\]
if the constant $c_0$ in \eqref{(4-18)} is small enough.
This proves \eqref{lastconde4}. 
\end{proof}

\end{document}
